\documentclass[12pt]{article}
\usepackage[utf8]{inputenc}
\usepackage[T1]{fontenc}
\usepackage[margin=1in]{geometry}
\usepackage{setspace}
\usepackage{xcolor}
\usepackage[colorlinks=true, urlcolor=blue!60!black, linkcolor=blue!60!black]{hyperref}
\setlength{\parindent}{2em}
\setlength{\parskip}{1.25em}
\renewcommand{\baselinestretch}{1.0}
\begin{document}
\begin{center}
{\LARGE \textit{Romer v. Evans}, 517 U.S. 620 (1996)\par}
\vspace{0.8em}
{\large Case Brief --- Four-Page Edition\par}
\end{center}
\vspace{1em}

\section{Memory Jogger}

The Court struck down Colorado's Amendment 2 under the Equal Protection Clause, holding that a state constitutional ban on all state and local anti-discrimination protections for gay, lesbian, and bisexual people fails even rational-basis review because its sweeping disability is inexplicable by anything but animus toward the targeted class (\href{https://www.courtlistener.com/opinion/118027/romer-v-evans/}{517 U.S. 620}).

\section{Detailed Case Facts}

Before 1992, Aspen, Boulder, and Denver banned sexual-orientation discrimination in housing, employment, education, public accommodations, and services, and state measures---an executive order for state employees, an insurance rule, and college policies---did the same. In November 1992 Colorado voters approved Amendment 2, about 53\% to 47\%. Titled ``No Protected Status Based on Homosexual, Lesbian or Bisexual Orientation,'' it barred the State and all its subdivisions from adopting or enforcing any law or policy by which homosexual, lesbian, or bisexual orientation or conduct would entitle anyone to ``any minority status, quota preferences, protected status or claim of discrimination.''

Richard Evans, a gay Denver employee, other individuals, and the three cities sued Governor Roy Romer (who had opposed the measure), the Attorney General, and the State.

\section{Procedural History}

The Denver trial court preliminarily enjoined Amendment 2. In \textit{Evans I} (1993), the Colorado Supreme Court held that it was subject to strict scrutiny because it burdened a fundamental right to participate equally in the political process. On remand, the trial court found that gays and lesbians are not a suspect class but that the State's interests failed strict scrutiny, and it permanently enjoined the measure; \textit{Evans II} (1994) affirmed. The U.S. Supreme Court affirmed on May 20, 1996, on a different rationale: rational-basis review.

\section{Judicial Votes}

Affirmed, \textbf{6--3}.
\begin{itemize}
\item \textbf{Majority:} Justice Kennedy, joined by Justices Stevens, O'Connor, Souter, Ginsburg, and Breyer.
\item \textbf{Concurrences:} None.
\item \textbf{Dissent:} Justice Scalia, joined by Chief Justice Rehnquist and Justice Thomas.
\end{itemize}

\section{Holding}

Amendment 2 violates the Equal Protection Clause. Even a classification that burdens no fundamental right and targets no suspect class must bear a rational relationship to an independent and legitimate purpose. A law that identifies a single class by one trait and imposes a broad, undifferentiated disability---barring its members from seeking protection against discrimination at every level of government---raises an inference of animosity, and a ``bare \dots{} desire to harm a politically unpopular group'' is not a legitimate interest. The Court did not adopt the political-participation rationale, decide whether sexual orientation is a suspect class, or mention \textit{Bowers v. Hardwick}, leaving that tension for \textit{Lawrence v. Texas}.

\section{Analysis of Opinions}

\textbf{Majority (Kennedy).} Opening with Justice Harlan's \textit{Plessy} dissent, Kennedy rejected Colorado's claim that Amendment 2 merely denied ``special rights.'' It repealed local ordinances and state policies and barred every level of government from adopting new ones, imposing a ``special disability upon those persons alone'' and withholding safeguards ``taken for granted by most people.'' Ordinary rational-basis cases such as \textit{Williamson v. Lee Optical} involve narrow laws whose purposes can be discerned, but Amendment 2 ``fails, indeed defies, even this conventional inquiry.'' It is ``at once too narrow and too broad,'' and its breadth is so discontinuous with the asserted interests---landlords' and employers' associational freedom and conserving enforcement resources---that it raises ``the inevitable inference that the disadvantage imposed is born of animosity toward the class of persons affected.'' It is ``a classification of persons undertaken for its own sake.''

\textbf{Dissent (Scalia).} ``The Court has mistaken a \textit{Kulturkampf} for a fit of spite.'' Amendment 2 left general anti-discrimination laws intact and barred only preferential treatment; requiring a group to seek protection at the state rather than local level is no equal protection violation. If \textit{Bowers} lets a State criminalize homosexual conduct, it may rationally deny special protections. Coloradans could preserve traditional sexual mores by statewide referendum, as state constitutional polygamy bans show, and the Court sided with ``the lawyer class,'' acting from ``political will'' rather than judgment.

\section{Examples: Future Applications}

\begin{enumerate}
\item \textbf{Same-Side (Invalid): Ban on Protections for Unmarried Cohabitants.} A state constitutional amendment forbidding any government to protect people from discrimination based on ``unmarried cohabitation status'' imposes the same sweeping, single-trait disability.
\item \textbf{Same-Side (Invalid): Federal Nonrecognition of Same-Sex Marriages.} Stripping validly married same-sex couples of benefits across more than 1,000 federal laws imposes a broad status-based disability, as Kennedy later reasoned in \textit{United States v. Windsor}.
\item \textbf{Opposite-Side (Valid): Uniform Statewide Employment Code.} A statute preempting all local employment ordinances does not single out one group and rationally serves commercial uniformity.
\item \textbf{Opposite-Side (Valid): Narrow Fair-Housing Exemption.} Exempting owner-occupied duplexes and religious organizations from a sexual-orientation housing law is narrow and tied to associational and religious interests.
\item \textbf{Fence-Sitter: City Charter Amendment.} A single city's charter amendment repealing and forbidding local sexual-orientation protections resembles Amendment 2 but leaves state remedies open; the Sixth Circuit upheld Cincinnati's after \textit{Romer}.
\end{enumerate}

\section{Critique}

\textbf{Originalist and textualist critiques.} Critics fault the majority's failure to confront \textit{Bowers}---Lynn Baker's ``missing pages''---and its unannounced ``rational basis with bite,'' with no standard for separating legitimate moral legislation from animus.

\textbf{Progressive and doctrinal views.} Evan Gerstmann faulted the ``sloppy reasoning'' for leaving equal protection murkier. Others reconstructed it: Daniel Farber and Suzanna Sherry's pariah principle, Akhil Amar's bill-of-attainder analogy, and Cass Sunstein's praise of judicial minimalism.

\section{Key Quotations}

\begin{quote}
``A law declaring that in general it shall be more difficult for one group of citizens than for all others to seek aid from the government is itself a denial of equal protection of the laws in the most literal sense.''
\end{quote}
\noindent --- Justice Kennedy, Opinion of the Court, 517 U.S. at 633.

\begin{quote}
``We must conclude that Amendment 2 classifies homosexuals not to further a proper legislative end but to make them unequal to everyone else. This Colorado cannot do. A State cannot so deem a class of persons a stranger to its laws.''
\end{quote}
\noindent --- Justice Kennedy, Opinion of the Court, 517 U.S. at 635.

\begin{quote}
``[I]f the constitutional conception of `equal protection of the laws' means anything, it must at the very least mean that a bare \dots{} desire to harm a politically unpopular group cannot constitute a legitimate governmental interest.''
\end{quote}
\noindent --- Justice Kennedy, Opinion of the Court (quoting \textit{Department of Agriculture v. Moreno}), 517 U.S. at 634--35.

\end{document}
